% ====================================================================
% 第6章：工业实践与框架对比（LaTeX版本，含TikZ图）
% ====================================================================

\section{工业实践与框架对比}\label{ch:industry}

\subsection{本章概述}\label{sec:ch6-overview}

Agent Evolution 从论文原型走向工业系统时，核心问题不再只是``智能体能不能完成一次任务''，而是``能否在真实约束下稳定、可控、可观测、低成本地持续完成任务''。第~\ref{ch:evaluation}~章讨论了评估体系如何构成自进化的选择压力；本章进一步讨论工程生态：开发者实际使用哪些框架，开源项目反映出怎样的技术栈分布，生产部署暴露出哪些可靠性、成本和安全挑战，以及为什么许多 demo 很容易令人兴奋，却很难直接变成 production。

本章依据的本地素材包括 \texttt{raw-github/} 中 348 个 GitHub 仓库快照，以及 \texttt{raw-blogs/} 中 1306 个博客、视频、产品页和平台内容文件。GitHub 语料显示，工业生态已形成多个层次：一层是通用 Agent 框架，如 LangChain、LangGraph、CrewAI、AutoGPT、smolagents、Vercel AI SDK；一层是工作流与平台化系统，如 n8n、browser-use、Skyvern、Suna、Gumloop、Lindy、Langflow；一层是记忆、评估、MCP、工具生态和自进化 harness，如 Letta、Mnemosyne、AgentJet、Auto-Harness、A-Evolve、Darwinia、ClawCode。博客语料则显示，2026年5月的内容集中在教程、框架上手、生产部署、记忆、可观测性、成本优化和安全护栏：在 1304 条可解析记录中，\texttt{self\_evolving\_agent} 与 \texttt{tutorial\_or\_implementation} 是最大类别，YouTube、Product Hunt、Medium/TDS、知乎、CSDN、GitHub Blog/Docs 等构成主要传播渠道。

这说明 Agent 工业实践已进入``框架丰富、应用广泛、评估不足、生产门槛高''的阶段。开发者不缺 demo 框架，也不缺工具调用范式；真正稀缺的是可复现的可靠性证据、面向业务流程的评估数据、跨模型成本控制、权限与安全边界、持续监控与回滚机制，以及能够将失败样本转化为系统改进的闭环。

\subsection{主流 Agent 框架对比}\label{sec:ch6-frameworks}

\subsubsection{框架生态的分层：从链式编排到状态机、团队协作与平台化}

从 raw-github 快照看，Agent 框架大致可分为四类。第一类是``组件与编排层''，代表是 LangChain 与 LangGraph。LangChain 在快照中被描述为``agent engineering platform''，Stars 约 137k、Forks 约 22.7k、Commits 约 16,023，是语料中影响力最大的通用 LLM 应用与 agent 生态之一。LangGraph Stars 约 32.5k、Forks 约 5.5k、Commits 约 6,870，定位更明确：构建长运行、有状态、可恢复的 agent。两者共同代表了一种工程化路线：不把 agent 视作一次性 prompt，而是把工具、状态、分支、检查点、人工介入和部署运行时组合成可管理系统。

第二类是``高层多智能体协作层''，代表是 CrewAI、AutoGen、MetaGPT 等。CrewAI 在 raw-github 中有官方快照，Stars 约 51.8k、Forks 约 7.2k、Commits 约 2,422，README 强调 role-playing autonomous agents、Crews 与 Flows，并将 Flows 描述为面向 enterprise/production 的事件驱动架构。AutoGen 的官方仓库未出现在本次 raw-github 快照中，但 raw-blogs 中有 23 条 AutoGen 相关教程或讨论，说明其在教学和多智能体对话范式中仍具存在感。此类框架的优势是降低多 agent 任务拆分、角色分工、handoff 和协作编排的门槛；弱点是抽象层较厚，真实业务中很容易出现``角色很多、调用很多、状态混乱、成本失控''的问题。

第三类是``代码/工具优先的轻量 agent 层''，代表是 Hugging Face smolagents、Vercel AI SDK、OpenAI Agents SDK 相关生态、PydanticAI、Agno 等。smolagents 在快照中 Stars 约 27.4k，强调 agents that think in code；Vercel AI SDK Stars 约 24.4k，面向 TypeScript 应用和 Web 产品集成。Agno 官方仓库未在 raw-github 中出现，但 raw-blogs 有 11 条 Agno 相关内容，主要围绕``15分钟构建 agent''``股票研究多智能体''``agent memory''等教程。此类框架通常强调少量代码、结构化输出、工具接口、模型提供商适配和前端/后端集成，适合从应用开发切入，而不一定覆盖完整生命周期。

第四类是``平台与工作流层''，代表是 AutoGPT、n8n、Langflow、Suna、Skyvern、Gumloop、Lindy 等。AutoGPT Stars 约 184k、Forks 约 46.2k，仍是从 2023 年 autonomous agent 热潮延伸出的标志性项目；n8n Stars 约 189k、Forks 约 57.8k，虽然本质是 workflow automation，但其 native AI capabilities 说明工作流平台正在吸收 agent 能力。Product Hunt 与 YouTube 语料中，Suna、Skyvern、Gumloop、Lindy、Pulze、Knowly 等产品频繁出现，反映出市场更关心``能否自动化浏览器、邮件、日程、销售、客服、QA、数据处理''等具体场景，而不是框架抽象本身。

图~\ref{fig:ch6-bubble}~以气泡图形式展示了主要框架在社区关注度（Stars）与工程成熟度两个维度上的相对位置。

% ====================================================================
% 图6.1: 框架对比气泡图（pgfplots）——Stars vs 工程成熟度
% ====================================================================
\begin{figure}[htbp]
\centering
\begin{tikzpicture}
\begin{axis}[
  width=14cm,
  height=9cm,
  xlabel={GitHub Stars (k)},
  ylabel={工程成熟度},
  xlabel style={font=\small},
  ylabel style={font=\small},
  xmin=0, xmax=200,
  ymin=0, ymax=6,
  ytick={1,2,3,4,5},
  yticklabels={L1 Demo, L2 Workflow, L3 Monitored, L4 Adaptive, L5 Evolving},
  yticklabel style={font=\scriptsize},
  xticklabel style={font=\scriptsize},
  grid=major,
  grid style={gray!20},
  title={主流 Agent 框架：Stars vs.\ 工程成熟度},
  title style={font=\small\bfseries},
  legend style={at={(0.02,0.98)}, anchor=north west, font=\tiny, draw=gray!50,
    fill=white, fill opacity=0.8, text opacity=1},
  clip=false,
]
% LangChain — 编排层
\addplot[only marks, mark=*, mark size=10pt, cat1!70, draw=cat1!90!black]
  coordinates {(137, 3)};
\addlegendentry{LangChain}
\node[font=\tiny, text=cat1!90!black, anchor=south] at (axis cs:137,3.25) {LangChain};

% LangGraph — 编排层
\addplot[only marks, mark=*, mark size=7pt, cat1!50, draw=cat1!80!black]
  coordinates {(32.5, 3.8)};
\addlegendentry{LangGraph}
\node[font=\tiny, text=cat1!80!black, anchor=south] at (axis cs:32.5,4.05) {LangGraph};

% CrewAI — 协作层
\addplot[only marks, mark=square*, mark size=8pt, cat2!70, draw=cat2!90!black]
  coordinates {(51.8, 2.8)};
\addlegendentry{CrewAI}
\node[font=\tiny, text=cat2!90!black, anchor=south] at (axis cs:51.8,3.05) {CrewAI};

% AutoGen — 协作层
\addplot[only marks, mark=square*, mark size=6pt, cat2!50, draw=cat2!80!black]
  coordinates {(25, 2.5)};
\addlegendentry{AutoGen}
\node[font=\tiny, text=cat2!80!black, anchor=south] at (axis cs:25,2.75) {AutoGen};

% Agno — 轻量层
\addplot[only marks, mark=triangle*, mark size=7pt, cat3!70, draw=cat3!90!black]
  coordinates {(12, 2.2)};
\addlegendentry{Agno}
\node[font=\tiny, text=cat3!90!black, anchor=south] at (axis cs:12,2.45) {Agno};

% smolagents — 轻量层
\addplot[only marks, mark=triangle*, mark size=7pt, cat3!50, draw=cat3!80!black]
  coordinates {(27.4, 2.6)};
\addlegendentry{smolagents}
\node[font=\tiny, text=cat3!80!black, anchor=north] at (axis cs:27.4,2.35) {smolagents};

% AutoGPT — 平台层
\addplot[only marks, mark=diamond*, mark size=8pt, cat4!70, draw=cat4!90!black]
  coordinates {(184, 1.8)};
\addlegendentry{AutoGPT}
\node[font=\tiny, text=cat4!90!black, anchor=south] at (axis cs:184,2.05) {AutoGPT};

% n8n — 平台层
\addplot[only marks, mark=diamond*, mark size=9pt, cat4!50, draw=cat4!80!black]
  coordinates {(189, 4.2)};
\addlegendentry{n8n}
\node[font=\tiny, text=cat4!80!black, anchor=south] at (axis cs:189,4.45) {n8n};

\end{axis}
\end{tikzpicture}
\caption{主流 Agent 框架气泡对比图。横轴为 GitHub Stars（千），纵轴为工程成熟度层级。颜色表示框架类别：蓝色=编排层，橙色=多智能体协作层，绿色=轻量agent层，紫色=平台与工作流层。气泡大小反映社区活跃度。}
\label{fig:ch6-bubble}
\end{figure}

\subsubsection{主要框架对比}

\begin{itemize}[nosep]
\item \textbf{LangChain}（raw-github: 137k stars, 16k+ commits）：通用 LLM 应用与 agent engineering platform。集成丰富、社区大、适配多模型/工具；但抽象演进快，项目需控制依赖复杂度。
\item \textbf{LangGraph}（raw-github: 32.5k stars, 6.8k+ commits）：有状态、长运行、可恢复 agent 编排。状态机、checkpoint、人类介入、生产运行时思路清晰；但学习曲线较高。
\item \textbf{CrewAI}（raw-github: 51.8k stars）：role-playing 多智能体与 Flows。上手快，角色/任务/团队抽象直观；但多 agent 容易增加 token 成本与调试难度。
\item \textbf{AutoGen}（raw-blogs 23 条相关教程）：多智能体对话、群聊、工具协作。适合教学与研究型 multi-agent prototype；但生产证据需另行验证，状态与权限边界需补强。
\item \textbf{Agno}（raw-blogs 11 条相关教程）：轻量 agent、memory、应用快速搭建。适合快速原型和垂直任务教程；但语料覆盖较少，社区规模与长期维护需观察。
\item \textbf{smolagents}（raw-github: 27.4k stars）：code-agent 与轻量实现。抽象小，便于理解和安全沙箱集成；但需开发者自行补齐复杂 orchestration。
\item \textbf{AutoGPT}（raw-github: 184k stars）：autonomous agent 平台与历史标志项目。社区关注度高，推动市场认知；但早期 autonomous loop 的可靠性争议仍是警示。
\item \textbf{n8n / Langflow}（raw-github 或 Product Hunt/YouTube 高频出现）：工作流/低代码 agent 编排。对业务用户友好，容易接入企业流程；但复杂推理与长期记忆能力依赖外部 agent 层。
\end{itemize}

框架对比不能只看 stars。LangChain/AutoGPT/n8n 的 star 数反映传播与历史积累，LangGraph/CrewAI 反映 agent 工程抽象的深化，smolagents/Vercel AI SDK/Agno 代表更轻量的开发者体验，Skyvern/Suna/Gumloop/Lindy 则代表产品化任务自动化。工业选型应首先回答四个问题：任务是否需要长状态？是否需要多 agent 协作？是否要接入企业权限与审计？是否有足够评估数据支撑自动化？不同答案会导向完全不同的框架组合。

\subsubsection{LangChain/LangGraph 与 CrewAI/AutoGen/Agno 的路线差异}

LangChain/LangGraph 的路线更像``把 agent 当作软件系统''。它重视组件、状态、图、checkpoint、运行时和监控，适合需要长期运行、可恢复、可审计的复杂任务。其不足在于工程门槛较高：开发者必须显式定义状态、边、条件、工具结果和异常恢复，否则图结构会变成另一种形式的复杂 spaghetti code。

CrewAI/AutoGen 的路线更像``把 agent 当作团队''。它们将任务拆成角色、职责、对话、handoff 和 manager-worker 关系，符合人类协作直觉，适合需求分析、报告生成、销售研究、内容流水线等工作。但团队隐喻也会掩盖底层事实：每个 agent 本质上仍是昂贵且不稳定的模型调用，角色越多，错误传播、上下文漂移和成本膨胀越严重。生产中必须用单元任务评估、权限隔离、固定输出 schema 和执行预算约束多 agent 系统。

Agno、smolagents、Vercel AI SDK 等路线则更强调``把 agent 当作应用组件''。它们适合直接嵌入 Web、数据分析、自动化脚本或轻量服务。优势是开发速度快、依赖少、容易与现有应用栈融合；风险是生命周期能力不足，如长期记忆、复杂回滚、跨会话评估和企业安全通常需要自行补齐。

\subsection{生产环境部署挑战}\label{sec:ch6-production}

\subsubsection{可靠性天花板：从``偶尔成功''到``持续可用''}

Agent demo 最常见的幻觉，是把一次成功轨迹误认为稳定能力。真实生产环境要求的是长期成功率、错误恢复、边界可控和可解释失败。raw-blogs 中``production''``reliable''``observability''``human-in-the-loop''等主题虽数量不及教程类内容，但它们指向同一问题：企业不缺能跑通 demo 的 agent，缺的是在失败时可诊断、可暂停、可回滚的 agent。

可靠性天花板来自多层叠加。模型层会出现幻觉、格式漂移、长上下文遗忘和工具选择错误；编排层会出现状态不一致、循环失控、handoff 失败和并发竞争；工具层会受到 API 变更、网页结构变化、权限不足和网络不稳定影响；数据层会出现检索污染、记忆过期和隐私泄漏；评估层则常常没有足够真实任务覆盖所有异常路径。任何一层失败，都可能让 agent 从``智能自动化''退化为``昂贵随机过程''。

因此，生产 agent 的基本工程模式应包括：任务拆分与最小可验证步骤、结构化输出与 schema 校验、工具调用白名单、超时与预算限制、状态 checkpoint、自动重试与降级策略、人类审批节点、失败样本归档、回归测试和线上监控。LangGraph 的状态化路线、CrewAI Flows 的事件驱动控制、OpenAI Agents SDK 相关教程中的 guardrails、Camunda/BPMN 中 human-in-the-loop 的流程思想，都在不同层面回应这个问题。

\subsubsection{成本挑战：token、工具、评估与人类审查共同构成总成本}

Agent 成本不只是 LLM token。一个生产 agent 往往包含规划调用、工具选择调用、执行调用、反思调用、评估调用、日志与监控、外部 API、浏览器环境、向量数据库、沙箱计算、人类审核和失败重跑。多 agent 框架会进一步放大成本：一个 manager 分配任务，多个 worker 执行，再由 reviewer 审查，表面上接近人类团队，实际上可能把一次业务操作扩展成几十次模型调用。

raw-blogs 中有关 LLM cost optimization、multi-agent 成本、agent workflow optimization 的内容虽然抓取质量不一，但足以说明成本已成为实践者关心的主题。工业系统需要的不只是``平均一次任务多少钱''，还要记录 cost per success、cost per accepted output、cost per regression prevented、cost per human review saved。否则，一个 agent 在 benchmark 上提升 5\%，却把线上成本提高 3 倍，很可能并不值得部署。

成本控制的常见策略包括：用规则和小模型处理低风险步骤；只在不确定性高的节点调用大模型；缓存工具结果和中间推理；将长上下文压缩为可验证状态；限制 agent 的最大步数和最大并发；用离线评估筛选 prompt、tool policy 和模型组合；把失败样本转化为测试集，避免同类错误反复消耗 token。对自进化 agent 而言，成本还必须进入 fitness function：进化目标不能只优化成功率，也要优化单位成本、单位延迟和单位风险。

\subsubsection{安全挑战：权限、数据、工具与自修改边界}

Agent 的安全风险高于普通聊天机器人，因为它不仅生成文本，还会调用工具、读写文件、访问网页、执行代码、触达企业系统，甚至在自进化场景中修改自身 prompt、记忆或代码。raw-github 中多个项目包含 \texttt{.env.example}、Dockerfile、\texttt{SECURITY.md}、\texttt{AGENTS.md}、CI workflow 等文件，说明开源生态已经意识到运行环境和贡献流程的重要性；但这并不等于默认安全。

安全边界至少包括四类。第一是权限边界：agent 能访问哪些 API、文件、数据库、浏览器页面和第三方工具，必须由最小权限原则限制。第二是数据边界：企业知识库、用户隐私、日志、记忆和训练样本不能被任意注入 prompt 或泄露到外部模型。第三是执行边界：代码执行、浏览器自动化和 shell 命令必须使用沙箱、审计、超时和网络限制。第四是自修改边界：如果 agent 可以修改 prompt、工具配置或代码库，就必须有不可被 agent 绕过的评估器、回滚机制和人工审批。

从 demo 到 production，安全不是最后加的插件，而是系统架构的一部分。尤其在 Agent Evolution 场景中，系统会主动寻找更高分策略；如果评估器只奖励任务完成，不惩罚越权、隐私泄露、测试绕过和不可维护代码，那么 agent 很容易出现 Goodhart 式优化。工业实践必须把安全事件、权限拒绝、策略违规、异常工具调用和人工否决率纳入与成功率同等重要的指标。

\subsubsection{可观测性与运维：从黑盒轨迹到可审计系统}

Agent 可观测性比传统服务更复杂，因为失败原因可能出现在 prompt、模型采样、检索结果、工具返回、状态转换、记忆选择、人类输入和外部环境中。raw-blogs 中出现 W\&B agent observability、Langfuse、LangGraph production、Docker deployment 等主题，反映实践者已从``如何构建 agent''转向``如何理解 agent 为什么这样做''。

生产系统至少需要记录：输入、任务分解、模型版本、prompt 版本、工具列表、每次调用参数、工具返回、状态变化、token 成本、延迟、异常、人工审批、最终结果和评价标签。更进一步，还需要 trace diff：同一任务在不同模型、不同 prompt、不同工具策略下为何表现不同。没有这些日志，自进化闭环无法可靠运行，因为系统无法判断一次失败应该归因于模型能力、工具错误、评估不充分，还是 prompt/状态设计缺陷。

\subsection{开源生态分析}\label{sec:ch6-opensource}

\subsubsection{Stars 与社区活跃度：头部项目吸收注意力，长尾项目快速实验}

raw-github 的 348 个仓库全部包含 star 信息。头部项目呈现明显平台化特征：n8n 约 189k stars、AutoGPT 约 184k、LangChain 约 137k、browser-use 约 94.8k、awesome-mcp-servers 约 87.2k、modelcontextprotocol/servers 约 86k、Zed 约 83.3k、CrewAI 约 51.8k、DSPy 约 34.5k、LangGraph 约 32.5k。这个排序说明，开发者关注的不只是``agent 算法''，还包括工作流自动化、浏览器操作、MCP 工具生态、开发环境、提示/程序化语言模型框架和多智能体框架。

但 star 不是生产成熟度。AutoGPT 的高 star 很大程度来自历史热潮和大众传播；n8n 的高 star 来自成熟工作流自动化与 AI 能力融合；LangChain 的高 star 来自广泛集成；LangGraph 的相对较低 star 反而可能代表更专业的生产编排需求。长尾项目如 A-Evolve、Auto-Harness、AgentJet、Darwinia、Mnemosyne、ClawCode stars 较少，却更接近 Agent Evolution 的前沿。

因此，开源生态可概括为``头部平台吸收入口流量，长尾项目探索能力边界''。工业用户选型时，应把 stars 作为社区风险指标，而不是能力指标。

\subsubsection{技术栈分布：Python 主导 agent 逻辑，JS/TS 连接产品界面}

对 348 个 raw-github 快照进行文件级启发式统计，Python 相关信号出现在约 207 个仓库，JS/TS 相关信号出现在约 42 个仓库，Docker 约 50 个，Rust 约 17 个，Go 约 2 个；同时，docs/test/CI 信号非常普遍，分别出现在 300+ 个仓库级别。图~\ref{fig:ch6-techstack}~展示了这一分布。

% ====================================================================
% 图6.2: 技术栈分布柱状图
% ====================================================================
\begin{figure}[htbp]
\centering
\begin{tikzpicture}
\begin{axis}[
  ybar,
  bar width=18pt,
  width=12cm,
  height=7cm,
  xlabel={技术栈},
  xlabel style={font=\small},
  ylabel={仓库数量},
  ylabel style={font=\small},
  ymin=0, ymax=340,
  symbolic x coords={Python, {JS/TS}, Docker, {Docs/CI}, Rust, Go},
  xtick=data,
  xticklabel style={font=\small},
  yticklabel style={font=\scriptsize},
  nodes near coords,
  nodes near coords style={font=\small\bfseries},
  title={Agent 开源项目技术栈分布（N=348 仓库）},
  title style={font=\small\bfseries},
  grid=major,
  grid style={gray!20},
  enlarge x limits=0.15,
]
\addplot[fill=cat1!60, draw=cat1!80!black] coordinates {
  (Python, 207)
  ({JS/TS}, 42)
  (Docker, 50)
  ({Docs/CI}, 300)
  (Rust, 17)
  (Go, 2)
};
\end{axis}
\end{tikzpicture}
\caption{Agent 开源项目技术栈分布。Python主导agent逻辑层，JS/TS连接产品界面，Docs/CI信号覆盖几乎所有项目。数据基于raw-github 348个仓库快照。}
\label{fig:ch6-techstack}
\end{figure}

这也解释了为什么 Vercel AI SDK、LangGraph.js、n8n、Langflow、browser-use web-ui 等项目重要。企业落地不是把 Python notebook 放到线上，而是把 agent 融入已有产品、权限、界面、消息系统、CRM、工单、数据平台和 CI/CD。Python 负责 agent 脑和评估闭环，JS/TS 负责用户触点和产品化体验，Docker/CI 负责可复现部署，观测平台负责闭环学习。

\subsubsection{时间切片与内容热度：2026年5月集中爆发，但存在采样偏差}

raw-github 的 time\_slice 中，2026-05 占 186 个，unknown 占 109 个；raw-blogs 的记录则高度集中在 2026-05。这表明本项目语料更像一个``2026年5月 Agent Evolution/Agent 工业生态快照''，而不是多年纵向数据集。因此，本章谈``star趋势''时采用的是相对热度和生态层级，而非严格增长曲线。若要进行真正的趋势分析，后续需要补充 GitHub star history、release history、issue/PR 活跃度、贡献者数量、下载量、Discord/论坛活跃度和企业采用案例。

即便如此，当前快照仍揭示了几个明确趋势。第一，MCP 和工具生态成为 agent 工业化的重要基础设施。第二，浏览器与工作流自动化是最直接的商业场景。第三，memory 和 stateful agent 成为从 demo 到 production 的关键差异。第四，自进化项目开始从论文转向 harness 与工程平台。

\subsubsection{社区活跃度的另一面：文档、测试、CI 与安全文件}

开源项目是否能支撑生产，不仅看 README，也看仓库结构。raw-github 中大量项目包含 \texttt{docs}、\texttt{tests}、\texttt{.github}、\texttt{Dockerfile}、\texttt{SECURITY.md}、\texttt{AGENTS.md}、\texttt{CLAUDE.md}、\texttt{pyproject.toml}、\texttt{package.json} 等文件。这说明生态已经从``prompt demo''逐步转向工程项目。

然而，文件存在不等于质量充分。工业团队应建立自己的 adoption checklist：许可证是否兼容、维护者是否活跃、release 是否稳定、issue 是否及时响应、是否有安全政策、是否支持离线/私有部署、是否能导出 trace、是否能限制工具权限、是否能接入内部评估和回滚。

\subsection{从 Demo 到 Production 的鸿沟}\label{sec:ch6-gap}

\subsubsection{Demo 成功依赖``窄任务 + 宽容环境''，生产成功依赖``宽任务 + 严格环境''}

Demo 通常选择窄任务、短流程、可控输入、低风险工具和人工挑选案例。生产环境恰好相反：输入不可控，业务规则复杂，用户会中断或纠错，权限系统繁琐，外部 API 会失败，网页会变化，数据有隐私要求，错误有真实成本。Agent demo 只需要展示一次``看起来像自主完成''；production 需要在数千次任务中稳定达成 SLA。

这就是从 demo 到 production 的第一道鸿沟：任务分布变化。图~\ref{fig:ch6-gap-diagram}~展示了这一鸿沟的关键屏障。

% ====================================================================
% 图6.3: Demo到Production鸿沟示意图（TikZ）
% ====================================================================
\begin{figure}[htbp]
\centering
\begin{tikzpicture}[
  >=Stealth,
  block/.style={
    rectangle, rounded corners=4pt, minimum width=3.2cm, minimum height=1.1cm,
    text centered, font=\small, draw=gray!60, line width=0.8pt,
    text width=3cm, align=center
  },
  barrier/.style={
    rectangle, rounded corners=3pt, minimum width=2.8cm, minimum height=0.7cm,
    text centered, font=\scriptsize, draw=gray!80, line width=0.6pt, fill=gray!10
  },
  arrow/.style={->, thick, gray!70},
]
% Left: Demo zone
\node[block, fill=cat3!25, draw=cat3!60] (demo) at (0, 0) {\textbf{Demo 成功区}\\窄任务 + 宽容环境};

% Five barriers
\node[barrier, fill=cat5!20, draw=cat5!50] (b1) at (5.5, 2.8) {可靠性天花板};
\node[barrier, fill=cat5!20, draw=cat5!50] (b2) at (5.5, 1.4) {成本失控};
\node[barrier, fill=cat5!20, draw=cat5!50] (b3) at (5.5, 0)   {安全边界缺失};
\node[barrier, fill=cat5!20, draw=cat5!50] (b4) at (5.5,-1.4) {可观测性不足};
\node[barrier, fill=cat5!20, draw=cat5!50] (b5) at (5.5,-2.8) {评估数据匮乏};

% Right: Production zone
\node[block, fill=cat1!25, draw=cat1!60] (prod) at (11, 0) {\textbf{Production 可用区}\\宽任务 + 严格环境};

% Gap label
\node[font=\footnotesize\bfseries, text=cat5!80!black, rotate=90] at (5.5, -4.0) {Demo $\to$ Production 鸿沟};

% Arrows from demo to barriers
\draw[arrow, cat3!60] (demo.east) -- (b1.west);
\draw[arrow, cat3!60] (demo.east) -- (b2.west);
\draw[arrow, cat3!60] (demo.east) -- (b3.west);
\draw[arrow, cat3!60] (demo.east) -- (b4.west);
\draw[arrow, cat3!60] (demo.east) -- (b5.west);

% Arrows from barriers to production
\draw[arrow, cat1!60] (b1.east) -- (prod.west);
\draw[arrow, cat1!60] (b2.east) -- (prod.west);
\draw[arrow, cat1!60] (b3.east) -- (prod.west);
\draw[arrow, cat1!60] (b4.east) -- (prod.west);
\draw[arrow, cat1!60] (b5.east) -- (prod.west);

% Annotations below each barrier
\node[font=\tiny, text=gray!60, anchor=north] at (5.5, 3.35) {幻觉/格式漂移/状态不一致};
\node[font=\tiny, text=gray!60, anchor=north] at (5.5, 1.95) {多轮调用放大 token 消耗};
\node[font=\tiny, text=gray!60, anchor=north] at (5.5, 0.55) {权限/数据/执行/自修改};
\node[font=\tiny, text=gray!60, anchor=north] at (5.5,-0.85) {缺少 trace/diff/审计日志};
\node[font=\tiny, text=gray!60, anchor=north] at (5.5,-2.25) {无真实任务失败样本};

\end{tikzpicture}
\caption{Demo 到 Production 的鸿沟示意图。五个关键屏障将Demo成功区与Production可用区隔开：可靠性天花板、成本失控、安全边界缺失、可观测性不足和评估数据匮乏。每一层屏障对应一类必须解决的工业落地挑战。}
\label{fig:ch6-gap-diagram}
\end{figure}

\subsubsection{工业落地需要``评估先行''，而不是``框架先行''}

许多团队的错误路径是先选框架，再找场景，最后才补评估。Agent Evolution 的经验恰好相反：应先定义任务集、成功标准、失败类别、成本上限和安全红线，再选择框架。若任务需要长期状态和人工审批，LangGraph/Camunda/BPMN 式流程可能更合适；若任务主要是内容流水线和角色协作，CrewAI/AutoGen 可快速验证；若任务是 Web 产品中的工具调用，Vercel AI SDK、OpenAI Agents SDK、PydanticAI、smolagents 可能更轻；若任务是业务流程自动化，n8n/Langflow/低代码平台可能更容易被组织接受。

评估先行还意味着建立``失败样本库''。每次线上失败都应归档为可重放任务，标注失败原因、影响范围、期望行为、修复策略和回归测试。自进化系统的价值不在于声称会自我改进，而在于能否把这些失败样本转化为 prompt、工具、记忆、代码、评估器和流程的可验证改进。如果没有这一闭环，agent 只是一次性自动化脚本；有了闭环，agent 才可能逐步接近工业级自治。

\subsubsection{Production 的成熟度模型}

结合 raw-github 与 raw-blogs 的生态信号，可以将 agent production maturity 分为五级。L0 是 prompt demo：只有提示词和一次性输出。L1 是 tool demo：能调用少量工具，但缺少状态和评估。L2 是 workflow agent：有固定流程、结构化输出、日志和人工审批。L3 是 monitored agent：有 trace、成本统计、错误分类、回归测试和安全护栏。L4 是 adaptive agent：能从失败样本中更新 prompt、工具策略、记忆或代码，并通过独立评估器验证收益。L5 才是真正的 self-evolving production agent：具备自动候选生成、隔离评估、成本/安全/性能多目标优化、灰度发布、自动回滚和长期审计。

当前开源生态多数项目处于 L1--L3。LangGraph、CrewAI Flows、OpenAI Agents SDK、n8n、Langflow 等帮助开发者到达 L2/L3；Auto-Harness、A-Evolve、AgentJet、DGM/SICA/AlphaEvolve 相关思想则指向 L4/L5。但 L4/L5 的门槛远高于框架安装：需要可执行评估器、真实任务数据、隔离 sandbox、版本管理、指标治理、组织流程和人类监督。

\subsubsection{本章结论：框架是入口，闭环才是壁垒}

工业实践与框架对比的结论可以概括为三点。第一，框架生态已经足够丰富，选型不应追求``唯一最佳框架''，而应基于任务状态、协作复杂度、产品栈、权限要求和评估能力组合不同工具。第二，从 demo 到 production 的主要瓶颈不是模型能否``想出下一步''，而是系统能否稳定记录、验证、限制、恢复和改进每一步。第三，Agent Evolution 的工业价值取决于闭环：失败是否被捕获，改进是否被验证，成本是否被控制，安全是否被强制，长期收益是否超过运维复杂度。

因此，未来的工业级 Agent Evolution 平台很可能不是单一 agent 框架，而是``框架 + 评估器 + 记忆 + 工具权限 + 可观测性 + 成本控制 + 灰度发布 + 自进化 harness''的组合系统。LangChain/LangGraph、CrewAI/AutoGen/Agno、smolagents/Vercel AI SDK、n8n/Langflow，以及 A-Evolve/Auto-Harness/AgentJet 等项目分别提供了这套系统的不同部件。真正的竞争不在于谁能写出更炫的 demo，而在于谁能把这些部件组织成可持续、可审计、可进化的生产基础设施。
